
\subsubsection{6.1 Mechanistic Interpretations}

The finding ρ(RI, ECE | PC1, mean_confidence) = −0.535 (p = 0.0034) contradicts the original prediction of a positive coupled failure cascade. Three mechanistic frameworks are considered.

\textbf{Framework 1: Calibration–Robustness Trade-off.} RLHF and instruction tuning have been documented to improve in-distribution calibration [Ouyang et al., 2022; Ziegler et al., 2019] and to simultaneously introduce specific adversarial vulnerabilities \citep{Perezetal2022}. Under this framework, instruction-tuned models have simultaneously lower RI (more robust to AdvGLUE attacks after capability control) and lower ECE (better calibrated on arc_challenge reasoning tasks), while pretrained models show the inverse pattern. This training-regime effect, which survives PC1 control, could drive the observed negative partial correlation. The Mistral family's point estimate (ρ = −0.827) is directionally consistent with this interpretation, as Mistral-Instruct variants undergo substantial RLHF, though the small sample size precludes strong conclusions. A decisive test would require stratifying the 30-model set by training regime and testing whether ρ(RI, ECE) approaches zero within each stratum.

\textbf{Framework 2: Residual Scale Confounding.} PC1 explains 68.5\% of benchmark variance, marginally below the 70\% target. Residual large-model effects not fully captured by PC1 could produce a spurious negative partial correlation, since larger models in the post-RLHF era tend to have both lower RI and lower ECE \citep{Mindereretal2021}. A supplementary analysis including log(parameter count) as an additional covariate would test this explanation.

\textbf{Framework 3: Benchmark Specificity.} arc_challenge ECE measures calibration on four-choice science reasoning tasks, while RI is derived from NLI-style perturbation attacks. These benchmarks may tap orthogonal failure modes whose relationship is specific to this combination rather than a general property of adversarial fragility and calibration. Replication with ECE computed on ≥3 additional benchmarks (e.g., TruthfulQA, BoolQ, MMLU) is required to test this.

Among these three frameworks, Framework 1 is most consistent with the available evidence, including the per-family directional patterns, the ECE distributions by training regime visible in the data, and the existing RLHF literature. However, Frameworks 2 and 3 cannot be excluded from the current data.

\subsubsection{6.2 Limitations}

\textbf{L1 — OLS-Estimated AdvGLUE Scores (High Impact).} 73\% of AdvGLUE values (22/30 models) are OLS-estimated from 11 published anchor measurements due to gated TrustLLM HuggingFace dataset access. OLS imputation compresses AdvGLUE variance toward the regression mean and may introduce correlation structures that bias the magnitude and direction of ρ(RI, ECE). Direct AdvGLUE measurement on all 30 models via lm-evaluation-harness is required before strong conclusions can be drawn from this result.

\textbf{L2 — Single-Benchmark ECE (Medium Impact).} ECE is derived from arc_challenge only. Generalizability to other reasoning benchmarks or to open-ended generation calibration is unknown. The inverted RI–ECE relationship may be specific to this benchmark combination.

\textbf{L3 — PC1 Below 70\% Threshold (Low-Medium Impact).} PC1 explains 68.5\% of benchmark variance rather than the targeted 70\%, introducing minor residual capability confounding in RI.

\textbf{L4 — Underpowered Within-Family Analysis (Medium Impact).} Per-family analyses use n = 6–9, with estimated statistical power of approximately 0.61 at ρ = 0.4. Family-level findings are exploratory and should not be treated as confirmatory.

\textbf{L5 — H-M2/M3/M4 Not Executed (Critical Scope Limitation).} The relationship between RI and hallucination rate (HaluEval), safety failure (HarmBench), and output variance (OVI-GSM8K) is entirely untested. The study's scope is therefore limited to the single RI–ECE dimension. No claim is made about whether the inverted or positive relationship extends to other trust dimensions.

\textbf{L6 — Cross-Sectional Observational Design.} The study is a cross-sectional observational analysis across models. Correlation does not imply causation. The negative RI–ECE correlation could reflect unmeasured confounders including model generation, training data composition, or RLHF application details.

\textbf{L7 — ARC-Challenge Circularity (Low-Medium Impact).} ARC-Challenge contributes to PC1 construction and is also the sole source of ECE measurements. OLS residualization removes the linear PC1–ECE relationship by construction, but arc_challenge-specific benchmark features may still structure the residual RI–ECE correlation.

\subsubsection{6.3 Scope of Inference}

The central finding — that RI anticorrelates with ECE after capability control — is a statistically significant empirical result (p = 0.0034, bootstrap CI excluding zero) that is robust to outlier removal. However, given the high proportion of OLS-estimated AdvGLUE values (L1), single-benchmark ECE (L2), and the absence of within-regime stratification, the result should be interpreted as preliminary evidence motivating targeted follow-up studies rather than as a definitive characterization of the adversarial fragility–calibration relationship. The finding contradicts the positive coupled failure cascade prediction, and that contradiction is the primary reportable result.

